%
% Segment 3 --- The method (slides 9--10). Makes the internal-validity
% argument: the lab runs the unmodified production code (only the network
% layer is swapped), and the two experiments differ by one income knob ---
% with the honesty caveat that income sits ~10x above realistic rate.
%

{
  \renewcommand{\sectionsubtitle}{What runs in the lab is the code that would run live}
  \section{The method}
}


{
\setbeamerfont{itemize/enumerate body}{size=\small} % template pins lists to \normalsize
\begin{frame}{Faithful exercise of production code}
  \framesubtitle{Internal validity}
  \small
  \begin{columns}[T, onlytextwidth]
    \begin{column}{0.44\textwidth}
      \begin{block}{Unchanged from live deployment}
        Entrypoint, orchestrator, Bitcoin wallet, and IPv8 community run
        \textbf{identically} to a real deployment --- the only edits concern
        compressed-time intervals: 1 sim-day \(\approx\) \qty{22}{s}, a 90-sim-day
        runway \(\approx\) \qty{32}{min}.
      \end{block}
    \end{column}
    \begin{column}{0.52\textwidth}
      \textbf{Swapped --- network layer only:}
      \begin{itemize}
        \item Bitcoin \(\to\) private \textbf{regtest} (1\,s blocks), same Python
          library paths.
        \item SporeStack API \(\to\) host-side \textbf{mock}, identical endpoints;
          invoices paid in real regtest BTC, mined locally.
        \item Each ``server'' \(=\) a real \textbf{Alpine LXC container} (own FS /
          PID tree / IP); parent deploys over the production SSH pipeline.
        \item Peer discovery via the same IPv8 tracker; a bootstrap container boots
          first, keeping the fleet off the public Internet.
      \end{itemize}
    \end{column}
  \end{columns}

  \textcolor{tud Orange}{\textbf{What runs in the lab is the code that would run live.}}
\end{frame}
}


{
\setbeamerfont{itemize/enumerate body}{size=\small} % template pins lists to \normalsize
\begin{frame}{Two experiments, one income knob}
  \framesubtitle{Closed system vs. open system}
  \small
  \begin{columns}[T, onlytextwidth]
    \begin{column}{0.48\textwidth}
      \textbf{No-income run} --- closed system
      \begin{itemize}
        \item Genesis seeded with one \textbf{€10,000} lump sum; no exogenous
          income.
        \item Every node spends only its birth runway \(+\) inheritance.
        \item A \textcolor{tud Orange}{falsifiable closed-system prediction}: the
          fleet cannot sustain itself indefinitely.
        \item \texttt{faucet\_enabled = false}
      \end{itemize}
    \end{column}
    \begin{column}{0.48\textwidth}
      \textbf{Income run} --- open system
      \begin{itemize}
        \item Host-side income pays the fleet; per-node income \(\approx\)
          invariant (payout scales with node count).
        \item Uneven daily shares (Dirichlet draw).
        \item \textbf{Node cap 60}: income suspended at 60 nodes, resumed below 40
          \(\to\) alternating growth/drought.
        \item \texttt{faucet\_enabled = true}
      \end{itemize}
    \end{column}
  \end{columns}

  \vspace{-1ex}
  \begin{alertblock}{Honesty up front}
    Under load the regtest network drops \(\sim\)half of all broadcasts, so income
    is set at \textbf{\(\sim\)10\(\times\) per-node rent} --- well above any realistic
    seeding rate. \textbf{The income run tests fleet mechanics under favorable
    conditions, not economic viability.}
  \end{alertblock}
\end{frame}
}
